% !TEX root = collapse_threshold_ru.tex
\documentclass[12pt,a4paper]{article}
\usepackage{cmap}
\usepackage[margin=1in]{geometry}
\usepackage{amsmath,amssymb,amsfonts}
\usepackage[utf8]{inputenc}
\usepackage[T2A]{fontenc}
\usepackage[russian]{babel}
\usepackage[hidelinks]{hyperref}
\usepackage{comfortaa}
\renewcommand*\familydefault{\sfdefault}
\usepackage[x11names]{xcolor}
\usepackage{indentfirst}

\makeatletter
\renewcommand{\thefootnote}{\arabic{footnote}}
\renewcommand{\@makefnmark}{\mbox{\textsuperscript{\textcolor{white}{\thefootnote}}}}
\renewcommand{\@makefntext}[1]{\parindent 1.8em\noindent\textsuperscript{\textcolor{white}{\thefootnote}}\ \textcolor{white}{#1}}
\renewcommand{\footnoterule}{\kern-3pt\color{white}\hrule width 0.4\textwidth height 0.4pt\kern 2.6pt}
\renewcommand{\ps@plain}{%
  \let\@oddhead\@empty
  \let\@evenhead\@empty
  \def\@oddfoot{\hfill\textcolor{white}{\thepage}\hfill}%
  \let\@evenfoot\@oddfoot
}
\makeatother

\pagestyle{plain}

\title{\textbf{Объективный коллапс ниже операционного порога}}
\author{%
  Олег Понфиленок\\[0.35em]%
  {\normalsize Независимый исследователь; \href{mailto:ponfil@gmail.com}{ponfil@gmail.com}}%
}
\date{}

\begin{document}

\pagecolor{black}
\color{white}

\maketitle

\begin{abstract}
Предлагается критерий объективного коллапса по \emph{всем} элементам приведённой матрицы плотности в базисе einselection: $|\rho_{ij}|<\varepsilon\Rightarrow\rho_{ij}=0$, где операционный пол $\varepsilon=1/N_{\mathrm{ost}}\sim 3\cdot 10^{-121}$ задан остатком вычислений causal patch, а $\Omega_{\max}=4\pi(R_{\mathrm{dS}}/\ell_P)^2\sim 10^{123}$~--- число адресов на горизонте де~Ситтера. Величина~$\varepsilon$~--- пол значений; пересечение пола~--- операционный порог объективности. Внедиагональ ниже~$\varepsilon$~--- необратимая декогеренция; диагональ ниже~$\varepsilon$~--- исход отсутствует в носителе. Порог агент-независим: элемент меньше одной оставшейся операции физически неразрешим внутри горизонта. Скремблирование~--- динамика достижения порога, не его определение; сложность распутывания, гашение эха Лошмидта и насыщение OTOC~--- операциональные корреляты. Два независимых маршрута~--- экспоненциальное затухание $|\rho_{01}|$ и рост $N_{\mathrm{scr}}(t)$~--- дают одно время достижения порога $t_{\mathrm{obj}}\sim 10^{-12}$--$10^{-11}\,\text{с}$ при комнатной температуре; длительность самого коллапса $t_{\mathrm{coll}}\sim\hbar/(k_B T)$ (десятки фемтосекунд при~$300\,\mathrm{K}$). Следствие для квантовых вычислений: потолок на неструктурированное скремблирование ($N_* \approx 400$); структурированные алгоритмы (Shor, QPE) остаются возможными.
\end{abstract}

\noindent\textbf{Ключевые слова:} объективная~декогеренция; операционный~порог; пол~$\varepsilon$; $\Omega_{\max}$; $N_{\mathrm{ost}}$; скремблирование; голографический~предел; эхо~Лошмидта; квантовые~вычисления

\section{Введение}

Стандартная квантовая механика допускает унитарную декогеренцию: приведённая энтропия подсистемы растёт, глобальная энтропия нулева, фаза записана в корреляциях с окружением и восстановима (эхо Лошмидта~\cite{swingle2018}). Классический бит при этом не рождён \emph{объективно}~--- интерпретация остаётся виртуальной.

Настоящая работа постулирует: в базисе einselection~\cite{zurek2009} все элементы приведённой плотности подчиняются одному операционному полу
\begin{equation}
|\rho_{ij}| < \varepsilon,\qquad \varepsilon=\frac{1}{N_{\mathrm{ost}}},
\end{equation}
ниже которого элемент онтологически равен нулю. Величина~$\varepsilon$~--- пол: ненулевой элемент не бывает мельче. Пересечение пола~--- порог объективности. Внедиагональ: декогеренция~--- не обратимое шифрование, а коллапс суперпозиции. Диагональ: исход с вероятностью меньше~$\varepsilon$ отсутствует в носителе. Здесь $\Omega_{\max}$~--- голографический бюджет адресов (разд.~3), $N_{\mathrm{ost}}$~--- остаток ортогональных операций causal patch (разд.~4), постулат~--- разд.~5. Скремблирование не входит в определение: это динамика, которой неструктурированная система достигает порога; эхо и OTOC~--- корреляты, не критерий. Связь с уилеровским «It from Bit»~\cite{wheeler1989}: из бита (объективной записи) возникает классическая реальность. Сначала оценим марковскую обратимость тепловой декогеренции (ревивалы); интуиция, ведущая к $\Omega_{\max}$, дана мысленным экспериментом далее.

\section{Марковская декогеренция, ревивалы и время обратимости}

Идея состоит в том, что объективный коллапс~--- это \emph{необратимая} декогеренция~\cite{bassi2013}, а декогеренция становится необратимой, когда запись среды перестаёт быть принципиально различимой внутри causal patch (разд.~5). До этого порога декогеренция в тепловой среде остаётся марковской и статистически обратимой~\cite{zurek2009}: фазы окружения перемешиваются за характерное тепловое время, а редкие флуктуации могут вернуть видимую когерентность (\emph{ревивал}; ср.\ эхо Лошмидта~\cite{swingle2018}). Оценим, сколько эффективных независимых степеней свободы ещё допускает естественный ревивал на космологическом горизонте наблюдения.

Максимально короткая тепловая шкала соответствует температуре Планка $T_P\approx 1{,}42\times 10^{32}\,\mathrm{K}$:
\begin{equation}
\tau_T \;=\; \frac{\hbar}{k_B T_P} \;=\; t_P \;\approx\; 5{,}39\times 10^{-44}\,\text{с}.
\end{equation}
Принимаем модель обычной хаотической тепловой среды: горизонт наблюдения~$10^{12}$~лет; минимально детектируемый возврат когерентности~$\varepsilon=0{,}01$; после каждого~$\tau_T$ фазы считаются статистически перемешанными; ревивал наступает, когда $N$ независимых фаз случайно попадают в область, дающую когерентность не меньше~$1\%$. Вероятность одной такой флуктуации
\begin{equation}
p(N,\varepsilon)\;\approx\;
\frac{\pi^{N/2}\,R^N}{\Gamma(N/2+1)\,(2\pi)^N},\qquad
R=\sqrt{2\ln(1/\varepsilon)},
\end{equation}
а среднее время ожидания
\begin{equation}
T_{\mathrm{rev}}\;\approx\;\frac{\tau_T}{p(N,\varepsilon)}.
\end{equation}
Численно: при $N=90$ получаем $T_{\mathrm{rev}}\approx 2{,}4\times 10^{11}$~лет; при $N=91$~--- $T_{\mathrm{rev}}\approx 1{,}9\times 10^{12}$~лет. Граница составляет примерно $N_{\mathrm{eff,max}}\approx 90$ эффективных независимых степеней свободы: тепловая хаотическая среда из $\sim 91$ независимо флуктуирующего эффективного кубита уже не даёт даже однопроцентного естественного ревивала в среднем за триллион лет~--- даже при максимально коротком планковском корреляционном времени.

Важно, что температура Планка~--- предел применимости обычного описания «тепловой среды». При энергиях порядка $k_B T_P\approx 1{,}22\times 10^{19}\,\text{ГэВ}$ нельзя надёжно использовать стандартную квантовую теорию поля без квантовой гравитации; такая плотность энергии сильно искривляет пространство-время и при локализации может вести к образованию чёрных дыр. Поэтому число~$90$~--- формальная экстраполяция выбранной статистической модели, а не установленный результат квантовой гравитации.

Для сравнения при том же горизонте и~$\varepsilon$: $1\,\text{мкК}$ даёт около~$44$ эффективных кубитов; $300\,\mathrm{K}$~--- около~$55$; температура Планка~--- около~$90$. Рост температуры на~$38$ порядков увеличивает границу лишь примерно вдвое, поскольку допустимое $N$ зависит от времени корреляции только логарифмически. Эти масштабы на много порядков ниже порога неструктурированного скремблирования $N_*\approx 400$ (разд.~6): марковские ревивалы ограничивают обратимость \emph{тепловой} записи, тогда как критерий $|\rho_{ij}|<1/N_{\mathrm{ost}}$ относится к онтологической неразрешимости элемента~$\rho$ внутри causal patch (разд.~5).

\section{Мысленный эксперимент: сфера направлений}

В центре~--- атом, запутанный с фотоном. Фотон уходит во все стороны сразу: сферическая суперпозиция направлений, без щелей. Вокруг~--- сферический экран радиуса~$R$ с датчиками по всей поверхности. Число различимых вариантов детекции~--- площадь сферы, делённая на площадь ячейки~$\lambda^2$ (как число микросостояний):
\begin{equation}
\Omega(R,\lambda) \;=\; \frac{4\pi R^2}{\lambda^2}.
\end{equation}
Каждая ячейка~--- один индивидуальный исход «фотон детектирован сюда».

Атом связан с квантовым компьютером, который \emph{кодирует направление}: в регистре записано, в какую из~$\Omega$ ячеек указывает исход. Чтобы задать все направления, нужно
\begin{equation}
N \;=\; \log_2 \Omega \;=\; \log_2(4\pi) + 2\log_2\!\left(\frac{R}{\lambda}\right)
\end{equation}
кубитов ($2^N$ состояний покрывают~$\Omega$ вариантов). Скремблирование~--- перемешивание этой записи направления в корреляциях «атом~+~регистр».

\textbf{Предельное скремблирование.} Максимальный радиус~--- горизонт де~Ситтера~$R_{\mathrm{dS}}$; минимальная длина волны~--- тепловая длина при температуре Планка~$T_P$:
\begin{equation}
R_{\max} \;=\; R_{\mathrm{dS}}, \qquad
\lambda_{\min} \;=\; \frac{\hbar c}{k_B T_P} \;=\; \ell_P,
\end{equation}
где $\ell_P=\sqrt{\hbar G/c^3}$~--- длина Планка. Тогда
\begin{equation}
\Omega_{\max} \;=\; \frac{4\pi R_{\mathrm{dS}}^2}{\lambda_{\min}^2} \;=\; 4\pi\left(\frac{R_{\mathrm{dS}}}{\ell_P}\right)^2 \;\sim\; 10^{123},
\end{equation}
\begin{equation}
N \;=\; \log_2 \Omega_{\max} \;=\; \log_2(4\pi) + 2\log_2\!\left(\frac{R_{\mathrm{dS}}}{\ell_P}\right) \;\approx\; 409.
\end{equation}
Индивидуальность исходов на сфере~--- именно~$\Omega$~--- задаёт голографический бюджет записи. Число $\Omega_{\max}$~--- одно: число адресов на горизонте; $1/\Omega_{\max}$~--- вес одной ячейки в максимально смешанном состоянии на этом пространстве. Ячейки сферы~--- базис, в котором живёт вся приведённая~$\rho$; рабочий операционный пол~$\varepsilon$ уточняется бюджетом операций (разд.~4--5). Геометрическое $N=\log_2\Omega_{\max}\approx 409$~--- сколько кубитов нужно, чтобы \emph{адресовать} все ячейки. Потолок неструктурированного скремблирования~$N_*$ (разд.~6) берётся не от адресов, а от порога вероятности: $N_*=\log_2 N_{\mathrm{ost}}\approx 400$, потому что оставшихся операций не хватает, чтобы различить все~$\Omega_{\max}$ исходов.

Это интуиция экспериментального принципа: если нельзя принципиально зафиксировать разницу в эксперименте, то разницы никакой нет. Поэтому более глубокое скремблирование не отличимо от предельного.

\textbf{Разрез «фотон\,$|$\,остальное».} Геометрическое $N\approx 409$~--- не потолок глобальной запутанности Вселенной и не предел на число частиц. Это максимум для \emph{одного} фотона~$A$ как подсистемы по числу ячеек: $\dim\mathcal{H}_A\le\Omega_{\max}=2^N$. Операционно различимых исходов не больше $N_{\mathrm{ost}}$, поэтому по разложению Шмидта энтропия запутанности с произвольной системой~$B$ (регистр, среда, «остальная Вселенная») удовлетворяет
\begin{equation}
S(A:B)\;\le\; N_*\;=\;\log_2 N_{\mathrm{ost}}\;\approx\; 400\ \text{эбит}.
\end{equation}
Запутать фотон с $500$ кубитами можно; лишние степени свободы~$B$ не создают дополнительных независимых каналов через границу~$A\,|\,B$. Ни один прибор в causal patch не выжимает из одного фотона больше~$N_*$ бит: адресация ячейки на горизонте упирается в остаток операций, а не в число планковских клеток. Регистр глубже~$N_*$ кубитов не кодирует больше различимых направлений.

\textbf{Масштабы.} Следует различать: $\Omega_{\max}\sim 10^{123}$~--- число адресов на сфере; $N=\log_2\Omega_{\max}\approx 409$~--- геометрические биты на адресацию ячеек; $N_{\mathrm{holo}}\sim 10^{122}$~--- порядок ёмкости всего causal patch (энтропия де~Ситтера в натах, разд.~4); $N_{\mathrm{ost}}\sim 10^{120}$~--- остаток ортогональных операций; $N_*=\log_2 N_{\mathrm{ost}}\approx 400$~--- операционный потолок канала и скремблирования. Тот же принцип $S_{\mathrm{ent}}(A:\mathrm{rest})\le\log_2 d_A$ применим к любой конечной подсистеме; число~$409$ здесь выведено для пространственной модели направления фотона, число~$400$~--- для порога вероятности.

\section{Вычислительный предел и операционный пол вероятности}

Сфера разд.~3 задаёт \emph{память} causal patch: $\Omega_{\max}$ адресов. Динамика космологического горизонта задаёт \emph{счёт}: сколько ортогональных операций материя внутри горизонта ещё может выполнить. Этот бюджет и есть операционный пол вероятности.

Возьмём горизонт Хаббла $R=c/H$ и второе уравнение Фридмана
\begin{equation}
\dot H = -\frac{4\pi G}{c^2}(\rho+p).
\end{equation}
Скорость роста горизонта:
\begin{equation}
\frac{dR}{dt}
= -c\,\frac{\dot H}{H^2}
= \frac{4\pi G}{c^3}(\rho+p)\,R^2.
\end{equation}
Излучением пренебрегаем ($\Omega_r<0{,}01\%$): $\rho=\rho_m+\rho_\Lambda$, $p_m=0$, $p_\Lambda=-\rho_\Lambda$, откуда $\rho+p=\rho_m$. Для плоской вселенной $\rho_m=\Omega_m\cdot 3c^4/(8\pi G R^2)$, и
\begin{equation}
\frac{dR}{dt} = \frac{3}{2}\,c\,\Omega_m.
\end{equation}

Энтропия горизонта $S=\pi k R^2/\ell_P^2$:
\begin{equation}
\frac{dS}{dt}
= \frac{2\pi k R}{\ell_P^2}\frac{dR}{dt}
= \frac{3\pi k c^4}{\hbar G}\,\Omega_m R.
\end{equation}
Энергия материи внутри горизонта $E_m=c^4\Omega_m R/(2G)$. Предел Ллойда~\cite{lloyd2000}
\begin{equation}
\dot N_m = \frac{2E_m}{\pi\hbar} = \frac{c^4\Omega_m R}{\pi\hbar G}.
\end{equation}
Тёмная энергия энтропию не производит ($\rho_\Lambda+p_\Lambda=0$). Отношение
\begin{equation}
\frac{\dot S}{k\,\dot N_m} = 3\pi^2
\end{equation}
--- $3\pi^2$~нат (около $42{,}7$~бит) на каждую ортогональную операцию материи. $\Omega_m$ и~$R$ сокращаются: связь голографическая и не зависит от эпохи, пока доминируют материя и~$\Lambda$.

Энтропия де~Ситтера конечна:
\begin{equation}
S_{\max}=\frac{\pi R_{\mathrm{dS}}^2}{\ell_P^2}\approx 3{,}3\times 10^{122}\ \text{нат}
\end{equation}
($\approx 4{,}8\times 10^{122}$~бит). Остаток вычислений
\begin{equation}
N_{\mathrm{ost}}
= \frac{S_{\max}-S}{3\pi^2}
= \frac{S_{\max}\Omega_m}{3\pi^2}
\approx 3{,}5\times 10^{120}.
\end{equation}
В терминах адресов сферы $N_{\mathrm{ost}}=\Omega_m\Omega_{\max}/(12\pi^2)$: операций меньше, чем ячеек, на фактор $12\pi^2/\Omega_m\sim 10^2$.

Любое измерение несёт информацию~\cite{landauer1961}. Бинарное «произошло\,/\,не произошло»~--- 1~бит; суммарное число таких попыток не превосходит $N_{\mathrm{ost}}$. Минимальная ненулевая вероятность, ещё совместимая с одним положительным исходом,
\begin{equation}
\varepsilon = \frac{1}{N_{\mathrm{ost}}}\approx 3\times 10^{-121}.
\end{equation}
Меньшая вероятность внутри causal patch неизмерима и операционно не существует. Отсюда же потолок неструктурированного канала: равномерный вес $1/2^{N}$ падает ниже~$\varepsilon$ при
\begin{equation}
N_* = \log_2 N_{\mathrm{ost}} = \log_2(1/\varepsilon)\approx 400.
\end{equation}
$N_{\mathrm{ost}}$ убывает вместе с~$\Omega_m$; ниже берём \emph{текущее} значение (на лабораторных временах $\Omega_m$ заморожена). Геометрическое $\log_2\Omega_{\max}\approx 409$ считает ячейки; операционное $N_*\approx 400$~--- различимые исходы. Дефицитный ресурс~--- операции, не ячейки.

\section{Порог объективности}

\paragraph{Постулат.}
В базисе einselection~\cite{zurek2009} (ячейки сферы, разд.~3) все элементы приведённой матрицы плотности подчиняются одному полу разд.~4:
\begin{equation}
|\rho_{ij}| < \varepsilon \;\Rightarrow\; \rho_{ij}:=0
\quad\text{для любых }i,j,\qquad
\varepsilon=\frac{1}{N_{\mathrm{ost}}}\approx 3\times 10^{-121}.
\end{equation}
Затем перенормировка $\mathrm{tr}\,\rho=1$. Если обнуляется диагональ~$i$, обнуляются строка и столбец~$i$ (иначе $\rho\not\ge 0$). Пол~$\varepsilon$ задаёт допустимые значения; порог объективности~--- момент, когда элемент его пересекает. Скремблирование~--- типичная динамика пересечения, не определение порога. Базис фиксирует einselection: критерий~--- свойство состояния относительно указательного базиса среды, а не произвольного разложения.

Два режима одной карты. Положительность даёт $|\rho_{ij}|^2\le\rho_{ii}\rho_{jj}$. Пока указательные исходы живы ($p_i\sim O(1)$), порог первой пересекает \emph{внедиагональ}: когерентность меньше~$\varepsilon$ физически неразрешима внутри горизонта; фазовый регистр аннигилирован; коллапс суперпозиции объективен. \emph{Диагональ} ниже~$\varepsilon$: исход отсутствует в носителе. Для типичного кубита в тепловой бане диагонали остаются $\sim 1/2$; $t_{\mathrm{obj}}$ по-прежнему задаётся затуханием $|\rho_{01}|$.

\textbf{До порога} ($|\rho_{ij}|\ge\varepsilon$): элемент ещё различим как доля вычислительного бюджета; глобальная энтропия может оставаться нулевой, рост $S(\rho)$ подсистемы~--- убыль во внедиагоналях, не онтологический акт.

\textbf{На пороге} ($|\rho_{ij}|<\varepsilon$): элемент меньше одной оставшейся операции causal patch~--- не «трудно записать», а негде актуализировать различие. Геометрический минимум одного адреса $1/\Omega_{\max}\sim 10^{-123}$ уже не достигается: операций не хватает раньше, чем ячеек.

Почему не сложность распутывания как определение. Сложность $C$ схемы, восстанавливающей когерентность, агентно окрашена (чьи ресурсы?) и при линейном (или полиномиальном) росте $C(t)$ достигает величины $\sim\Omega_{\max}$ лишь за космологические времена,~--- тогда как объективация должна происходить быстро. «Число операций» прямой схемы наследует ту же линейность. Критерий по $|\rho_{ij}|$ агент-независим и при экспоненциальном затухании внедиагонали пересекается за пикосекунды (см.\ ниже).

Операциональные корреляты~--- не определение: необратимое гашение эха Лошмидта~\cite{swingle2018}; насыщение OTOC (\emph{out-of-time-order correlator}~--- вневременной коррелятор)~\cite{maldacena2016} на неструктурированной динамике; рост сложности распутывания как зонд той же динамики.

Следствие для носителя. Тот же пол на диагонали обрезает носитель состояния. Интегральный хвост кулоновского $1s$ уже на нанометровом удалении от ядра несёт вес ниже~$\varepsilon$: конфликта со спектроскопией нет, но носитель волновой функции компактен~--- это модификация уравнения Шрёдингера, а не только декогеренция.

\paragraph{Величина порога.}
Из разд.~3--4:
\begin{equation}
\ln\Omega_{\max}\approx 283,\qquad
\ln N_{\mathrm{ost}}\approx 278,\qquad
\varepsilon\sim 3\times 10^{-121}.
\end{equation}
$\Omega_{\max}$ совпадает по порядку с числом битов на горизонте де~Ситтера $N_{\mathrm{holo}}$ (ср.\ Кароляхази~\cite{karolyhazy1966}). Рабочий пол~--- $1/N_{\mathrm{ost}}$; потолок скремблирования~--- $N_*=\log_2 N_{\mathrm{ost}}\approx 400$, а не $\log_2\Omega_{\max}\approx 409$. Отождествление «сложность${}={}$адреса${}={}$ёмкость${}={}$операции» снимается.

\paragraph{Два времени: достижение порога и сам коллапс.}
Постулируя экспоненциальное затухание $|\rho_{01}|\sim e^{-t/\tau_{\mathrm{th}}}$ с тепловым временем $\tau_{\mathrm{th}}=\hbar/(k_B T)$, следует различать две шкалы.

\emph{Время достижения порога}~$t_{\mathrm{obj}}$~--- сколько ждать, пока внедиагональ упадёт ниже~$\varepsilon$. Условие $|\rho_{01}|<e^{-\ln N_{\mathrm{ost}}}$ даёт
\begin{equation}
t_{\mathrm{obj}} = \ln N_{\mathrm{ost}}\,\tau_{\mathrm{th}}\approx 278\,\tau_{\mathrm{th}}.
\end{equation}
Геометрический маршрут $|\rho_{01}|<1/\Omega_{\max}$ даёт $283\,\tau_{\mathrm{th}}$; на заявленной точности это те же пикосекунды. Это полное время накопления декогеренции до критерия, а не длительность акта коллапса: почти всё это время система ещё \emph{выше} порога.

\emph{Длительность самого коллапса}~$t_{\mathrm{coll}}$~--- время пересечения окрестности порога, когда онтологический акт уже происходит. При экспоненциальном законе окрестность $|\rho_{01}|\sim\varepsilon$ проходится за одну тепловую шкалу:
\begin{equation}
t_{\mathrm{coll}} \sim \tau_{\mathrm{th}} = \frac{\hbar}{k_B T}.
\end{equation}
При $300\,\mathrm{K}$ это $\sim 2{,}5\times 10^{-14}\,\text{с}$ (десятки фемтосекунд); по соотношению неопределённостей энергия этой шкалы $\hbar/t_{\mathrm{coll}}\sim k_B T$~--- температура бани, а не $\sim 1\,\mathrm{K}$, отвечающая всей $t_{\mathrm{obj}}$.

Численно (столбец $t_{\mathrm{obj}}$~--- по $283\,\tau_{\mathrm{th}}$; при $278\,\tau_{\mathrm{th}}$ те же порядки):
\begin{center}
\begin{tabular}{lcc}
\hline
$T$ & $\tau_{\mathrm{th}}=t_{\mathrm{coll}}$ & $t_{\mathrm{obj}}\approx 283\,\tau_{\mathrm{th}}$ \\
\hline
$300\,\mathrm{K}$ & $2{,}55\times 10^{-14}\,\text{с}$ & $7{,}2\times 10^{-12}\,\text{с}$ \\
$77\,\mathrm{K}$ & $9{,}9\times 10^{-14}\,\text{с}$ & $2{,}8\times 10^{-11}\,\text{с}$ \\
$4\,\mathrm{K}$ & $1{,}9\times 10^{-12}\,\text{с}$ & $5{,}4\times 10^{-10}\,\text{с}$ \\
\hline
\end{tabular}
\end{center}

\paragraph{Консилиенция.}
Независимый маршрут~\cite{entropy_threshold2025}: рост $N_{\mathrm{scr}}(t)\sim (t/\tau_{\mathrm{th}})^3$ до $N_*=\log_2 N_{\mathrm{ost}}\approx 400$ даёт то же $t_{\mathrm{obj}}\sim 10^{-12}$--$10^{-11}\,\text{с}$ при комнатной температуре. Два механизма~--- экспоненциальное гашение внедиагонали и объёмный рост скремблированной записи~--- сходятся к одному порядку \emph{времени достижения порога}. Эта консилиенция усиливает постулат: $\Omega_{\max}$ входит в геометрию адресов, $\varepsilon=1/N_{\mathrm{ost}}$~--- в операционный пол, $N_*$~--- в потолок неструктурированной динамики; длительность акта при этом остаётся $t_{\mathrm{coll}}\sim\tau_{\mathrm{th}}$.

Родственные пороговые критерии~\cite{entropy_threshold2025,oqt2026,bassi2013} обсуждают коллапс; здесь акцент~--- на агент-независимом пороге по всем $|\rho_{ij}|$ с геометрической привязкой к $\Omega_{\max}$ и вычислительной к $N_{\mathrm{ost}}$, а не на сложности распутывания как определении.

\section{Следствие для квантовых вычислений}

Порог бьёт по \emph{неструктурированному} скремблированию, а не по числу кубитов в регистре. Random circuit sampling загоняет релевантные внедиагонали записи среды ниже~$\varepsilon$ раньше, чем структурированные алгоритмы (Shor, QPE) при той же аппаратуре. Изоляция без смены логической структуры разрыв не закрывает: помогает структура цепи, удерживающая логическую суперпозицию вне хаотической записи среды.

Для полностью неструктурированных состояний эффективная ёмкость канала скремблирования ограничена порогом вероятности, а не числом ячеек:
\begin{equation}
N_* = \log_2 N_{\mathrm{ost}} = \log_2(1/\varepsilon)\approx 400.
\end{equation}
Это потолок Haar-подобной динамики и разреза «подсистема${}|$rest» (разд.~3--4), а не запрет на многокубитный структурированный регистр. Алгоритм Шора остаётся возможным: он ведёт логическую когерентность узким вычислительным путём и не отождествляется с заливкой амплитуд в голографический бюджет тепловой среды. Порог $\varepsilon\sim 10^{-121}$ относится к einselected-записи среды относительно оставшихся операций causal patch, а не к любой лабораторной амплитуде логического кубита под коррекцией ошибок.

Родственная идея~--- инвариантное множество Палмера~\cite{palmer2009}: реальность ограничена фрактальным инвариантным множеством; у Палмера объективность геометрическая и ближе к полу на всей~$\rho$, без операционального $\varepsilon=1/N_{\mathrm{ost}}$.

\section{Экспериментальная проверка}

Абсолютное значение $\varepsilon\sim 10^{-121}$ в лаборатории недостижимо как прямое измерение. Экспериментально проверяются \emph{операциональные корреляты} постулата (разд.~5): как при наращивании глубины скремблирования меняются восстановимость когерентности, эхо и OTOC и где унитарная картина перестаёт быть операционно адекватной.

\paragraph{Эхо Лошмидта.}
Система~+~среда подвергаются случайной унитарной динамике глубины~$d$ (слои двухкубитных вентилей, Haar-подобное распределение). Затем выполняется обратная эволюция $U^\dagger$ (эхо). \textbf{До порога}: видимость эха $F(d)$ восстанавливается при достаточной точности $U^\dagger$; рост $S(\rho)$ подсистемы~--- убыль во внедиагоналях. \textbf{На пороге}: $F(d)$ необратимо гаснет с $d$~\cite{swingle2018}; дополнительная глубина $U^\dagger$ не возвращает когерентность. Контроль: структурированная (не скремблирующая) цепь той же глубины~--- эхо сохраняется. Ось~--- $d$ и число кубитов~$N$; метрики~--- $F(d)$ и оценка $|\rho_{01}(d)|$ в указательном базисе.

\paragraph{OTOC (вневременной коррелятор).}
На неструктурированной динамике~\cite{maldacena2016} измеряют
\begin{equation}
F(t) = \frac{1}{Z}\operatorname{tr}\!\bigl(W(t)^\dagger V(0)^\dagger W(t) V(0)\bigr),
\end{equation}
где $W(t)=e^{iHt}We^{-iHt}$. До порога $F(t)$ отражает унитарное перемешивание; на пороге~--- насыщение и необратимое гашение чувствительности к ранним операторам при росте $t$ и~$N$. Скан по глубине скремблирования и размеру системы; сравнение со структурированными гамильтонианами (интегрируемыми, локально обратимыми).

\paragraph{Косвенный тест: структура vs.\ хаос.}
Разд.~6 предсказывает различие масштабов: random circuit sampling упирается в $N_* \approx 400$ раньше, чем структурированные алгоритмы (включая Shor) при той же аппаратуре. Полный $N_*$ недостижим, но сравнение \emph{глубины}, на которой классическая симуляция и эхо/OTOC ещё работают, для хаотических и структурированных схем одного $N$~--- дифференциальный тест: хаос должен «ломаться» раньше. Улучшение изоляции без смены логической структуры не должно сдвигать этот разрыв (разд.~6).

\paragraph{Космологический след.}
Ни согласование $\Omega_{\max}$ с $N_{\mathrm{holo}}$~\cite{karolyhazy1966}, ни вывод $N_{\mathrm{ost}}$ (разд.~4) не являются лабораторным тестом. Следствия~--- глобальный потолок на необратимое скремблирование записи среды и пол~$\varepsilon$~--- согласуются с голографической ёмкостью и остатком вычислений, но фальсифицируются только косвенно: через совокупность лабораторных коррелятов и астрофизических ограничений на нелинейную декогеренцию.

\bibliographystyle{unsrt}
\bibliography{refs}

\end{document}
